\documentclass[11pt,a4paper,oneside]{book}
\usepackage[T1]{fontenc}
\usepackage[utf8]{inputenc}
\usepackage{lmodern}
\usepackage{microtype}
\usepackage[a4paper,top=2.6cm,bottom=2.8cm,left=3.0cm,right=2.6cm]{geometry}
\usepackage{amsmath}
\usepackage{amssymb}
\usepackage{amsthm}
\usepackage{enumitem}
\usepackage{booktabs}
\usepackage{array}
\usepackage{xcolor}
\usepackage{graphicx}
\usepackage{tikz}
\usepackage{pgfplots}
\pgfplotsset{compat=1.18}
\usetikzlibrary{arrows.meta,positioning,calc,patterns}
\usepackage[font=small,labelfont=bf,justification=centering]{caption}
\usepackage[most]{tcolorbox}
\usepackage{makeidx}
\usepackage[hidelinks,bookmarksnumbered,bookmarksopenlevel=1,plainpages=false]{hyperref}

\hypersetup{
  pdftitle={Why Machines Can Be Wrong: From Linear Algebra to Machine-Learning Security (First Edition)},
  pdfauthor={Afiq Andico Pangimpian},
  pdfsubject={A tiered textbook: mathematics, computing, machine learning, and security},
  pdfkeywords={linear algebra, real analysis, probability, optimisation, neural networks, adversarial ML, Lean 4},
  pdfcreator={pdfLaTeX},
}

\definecolor{ink}{HTML}{1A1A1A}
\definecolor{accent}{HTML}{8C2F39}
\definecolor{layerorange}{HTML}{B45309}
\definecolor{layerteal}{HTML}{0F766E}
\definecolor{layergrey}{HTML}{374151}
\definecolor{layerindigo}{HTML}{4338CA}
\definecolor{layerjade}{HTML}{15803D}
\definecolor{layerbrown}{HTML}{92400E}
\definecolor{layerplum}{HTML}{86198F}
\color{ink}

% Environment names stay Indonesian so Indonesian and English chapters share one
% structural checker; only the printed labels are translated.
\theoremstyle{definition}
\newtheorem{definisi}{Definition}[chapter]
\theoremstyle{plain}
\newtheorem{teorema}{Theorem}[chapter]
\newtheorem{lemma}[teorema]{Lemma}
\newtheorem{akibat}[teorema]{Corollary}
\newtheorem{proposisi}[teorema]{Proposition}
\theoremstyle{remark}
\newtheorem{contoh}{Example}[chapter]
\newtheorem{latihan}{Exercise}[chapter]
\newtheorem{catatan}{Note}[chapter]
\newtheorem*{ringkasan}{Summary}

\renewcommand{\chaptername}{Chapter}
\renewcommand{\partname}{Part}
\renewcommand{\contentsname}{Contents}
\renewcommand{\bibname}{References}
\renewcommand{\indexname}{Index}
\renewcommand{\figurename}{Figure}
\renewcommand{\tablename}{Table}
\renewcommand{\listfigurename}{List of Figures}

\setlist[itemize]{topsep=0.35em,itemsep=0.15em,parsep=0pt}
\setlist[enumerate]{topsep=0.35em,itemsep=0.15em,parsep=0pt}

% ---------------- four reader layers ----------------
\newtcolorbox{untukanak}{colback=orange!6,colframe=layerorange,
  title={\small\bfseries For younger readers (age 10 and up)},
  fonttitle=\bfseries,breakable,left=6pt,right=6pt,top=4pt,bottom=4pt}
\newtcolorbox{untukdosen}{colback=teal!5,colframe=layerteal,
  title={\small\bfseries Note for instructors},
  fonttitle=\bfseries,breakable,left=6pt,right=6pt,top=4pt,bottom=4pt}
\newtcolorbox{untukpeneliti}{colback=black!4,colframe=layergrey,
  title={\small\bfseries For researchers: limits and next steps},
  fonttitle=\bfseries,breakable,left=6pt,right=6pt,top=4pt,bottom=4pt}
% ---- learning-style channels: auditory, kinesthetic, procedural, problem solving ----
\newtcolorbox{bacakeras}{colback=violet!5,colframe=layerindigo,
  title={\small\bfseries Read it aloud (auditory learner)},
  fonttitle=\bfseries,breakable,left=6pt,right=6pt,top=4pt,bottom=4pt}
\newtcolorbox{tangan}{colback=green!5,colframe=layerjade,
  title={\small\bfseries Do it by hand (kinesthetic learner)},
  fonttitle=\bfseries,breakable,left=6pt,right=6pt,top=4pt,bottom=4pt}
\newtcolorbox{resep}{colback=orange!4,colframe=layerbrown,
  title={\small\bfseries Step-by-step recipe (procedural learner)},
  fonttitle=\bfseries,breakable,left=6pt,right=6pt,top=4pt,bottom=4pt}
\newtcolorbox{tantangan}{colback=magenta!4,colframe=layerplum,
  title={\small\bfseries Challenge (problem solver)},
  fonttitle=\bfseries,breakable,left=6pt,right=6pt,top=4pt,bottom=4pt}

% In the English edition the first argument is the English term and the second is
% the Indonesian original, so the printed pair stays the same way round.
\newcommand{\kamus}[3]{%
  \par\smallskip\noindent
  \colorbox{blue!5}{\parbox{\dimexpr\linewidth-2\fboxsep}{%
    \small\textbf{#1} (\emph{#2}) --- #3}}\par\smallskip}
\newcommand{\tingkat}[1]{%
  \par\noindent{\footnotesize\scshape Level: #1}\par\smallskip}
\newcommand{\kertas}[1]{%
  \par\medskip\noindent
  \fbox{\parbox{\dimexpr\linewidth-2\fboxsep-2\fboxrule}{\small #1}}%
  \par\medskip}
\newcommand{\peta}[1]{%
  \par\medskip\noindent{\small\itshape Map to source texts: #1}\par\medskip}

\newenvironment{lab}{\par\medskip\noindent\textbf{Computational lab}\par\nopagebreak}{\par\medskip}
\makeatletter
\renewcommand{\verbatim@font}{\footnotesize\ttfamily}
\makeatother

\makeindex
\setlength{\emergencystretch}{3em}

% ---------------------------------------------------------------------------
% Running head untuk judul bab yang panjang.
%
% Kepala halaman hanya memuat "CHAPTER n. <judul>". Kalau judul penuh terlalu
% lebar, kepala halaman meluber (Overfull \hbox) dan itu bukan hal yang boleh
% dibiarkan. Menyingkat judulnya dilarang: halaman bab dan daftar isi harus
% memuat terjemahan utuh. Jadi yang dipendekkan hanya mark header, dan
% pemendekannya berlaku SEKLALI untuk bab berikutnya:
%
%   \judulpendek{Riemann Integrals}
%   \chapter{Riemann Integrals and the Fundamental Theorem of Calculus}
%
% Tanpa mekanisme sekali-pakai ini, \renewcommand{\chaptermark} di dalam satu
% berkas bab akan bocor ke semua bab sesudahnya - pernah terjadi: header bab
% 14-18 ikut berbunyi "CHAPTER 13. RIEMANN INTEGRALS".
\makeatletter
% \@chapapp (bukan \chaptername): setelah \appendix, \@chapapp menjadi "Appendix",
% sedangkan \chaptername tetap "Chapter", sehingga kepala halaman lampiran dulu
% berbunyi "CHAPTER A. TABLE OF NOTATION" tanpa satu pun peringatan LaTeX.
\newcommand*\judulpendek[1]{\gdef\kepala@pendek{#1}}
\newcommand*\kepala@pendek{}
\renewcommand*\chaptermark[1]{%
  \markright{\MakeUppercase{%
    \ifnum\c@secnumdepth>\m@ne\if@mainmatter\@chapapp\ \thechapter.\ \fi\fi
    \ifx\kepala@pendek\@empty #1\else\kepala@pendek\fi}}%
  \global\let\kepala@pendek\@empty}
\makeatother
